\documentclass[tikz,border=8pt]{standalone}
\usepackage{ctex}
\usepackage{tikz}
\usetikzlibrary{arrows.meta,positioning,fit,backgrounds}

\begin{document}
\begin{tikzpicture}[
  font=\small,
  >=Latex,
  blk/.style={draw=#1!72!black, fill=#1!8, rounded corners=2.2mm, very thick, minimum width=34mm, minimum height=11mm, align=center},
  flow/.style={->, very thick, draw=black!75}
]

\node[blk=blue] (can2d) at (0,1.6) {CAN 2D 输入\\\(64\times8\) 滑窗};
\node[blk=teal] (eth2d) at (0,-1.6) {Eth 2D 输入\\\(10\times80\) 序列};
\node[blk=purple] (align) at (5.5,0) {CrossDomainAligner\\对比学习 + MMD};
\node[blk=orange] (fuse) at (10.8,0) {融合时序\\\(\mathbf{Z}\in\mathbb{R}^{10\times d_f}\)};
\node[blk=green!70!black] (gt) at (16.1,0) {GraphTransformerIDS\\链类型 + 阶段输出};

\draw[flow] (can2d.east) -- (align.west);
\draw[flow] (eth2d.east) -- (align.west);
\draw[flow] (align.east) -- (fuse.west);
\draw[flow] (fuse.east) -- (gt.west);

\node[draw=black!35, fill=yellow!12, rounded corners=2mm, align=left, inner sep=4pt] at (10.8,-2.2) {
核心思想：先做跨域语义对齐，再做时间链建模；\\
输出同时包含链级结论与每时间步阶段分布。};

\begin{scope}[on background layer]
  \node[draw=black!25, rounded corners=3mm, line width=0.8pt,
  fit=(can2d)(eth2d)(align)(fuse)(gt), inner sep=8mm,
  label={[yshift=2mm]\large\bfseries CAN+以太网联合二维时序 ML 总览}] {};
\end{scope}
\end{tikzpicture}
\end{document}
